\section{国内外研究现状}\label{sec:related_work}

\subsection{通用目标检测方法演进}

通用目标检测的发展可以先从 R-CNN 系列看起。Girshick 等人的工作\cite{girshick2014_rcnn}把候选区域、特征提取和分类回归串成一条检测流程，后来 Faster R-CNN\cite{ren2015_fasterrcnn} 又把候选框生成进一步网络化，成为 two-stage 路线中很常被引用的代表。这一路线通常定位较稳，但候选区域处理会带来额外计算，推理速度不如更直接的单阶段方法。

与之相比，one-stage 方法取消了显式候选区域生成过程，直接在特征图上进行密集预测。YOLO 系列\cite{redmon2016_yolo,redmon2018_yolov3}、RetinaNet\cite{lin2017_focalloss} 和 EfficientDet\cite{tan2020_efficientdet} 均属于该路线的代表工作。它们结构紧凑、推理速度快、部署方便，因此在工程场景中应用较多；但遇到复杂背景和弱目标时，误检与漏检问题仍需要单独分析。

anchor-free 方法的出发点是减少 anchor 尺寸、比例和匹配策略带来的调参负担。FCOS\cite{tian2019_fcos} 直接在特征图位置上回归目标框，并用 centerness 抑制偏离目标中心的预测。这个设计在通用场景中较有代表性，但遥感微小船舶的目标面积太小，anchor-free 方法能否稳定优于 anchor-based 或专项方法，不能只依据通用数据集经验判断，还需要放到同一任务口径下测试。

另一条近年的路线是 transformer 检测器。DETR\cite{carion2020_detr} 将目标检测组织成端到端集合预测问题，Deformable DETR\cite{zhu2021_deformabledetr} 和 DINO\cite{zhang2023_dino} 继续改进收敛速度和检测精度。这类方法的全局建模能力较强，不过训练资源、数据规模和调参时间要求也更高。考虑到本课题需要同时完成实验比较和系统实现，本文将其作为方法脉络的一部分讨论，没有把它列为主要复现对象。

把这些方法放到一起看，差异不只体现在精度上，还体现在速度、输入设置、调参成本和系统接入方式上。本文选择 DRENet、FCOS 与 YOLO26，就是为了保留专项方法、anchor-free 方法和 one-stage 方法三个参照面。

\subsection{小目标检测研究进展}

小目标检测之所以困难，首先是因为可用像素少。目标在特征图上经过下采样后可能只剩很小响应，纹理、边缘和上下文都不稳定；同时，图像中大量背景区域会放大正负样本不平衡。围绕这些问题，相关研究常从下面几类思路展开。

一类做法直接改进特征表达。FPN\cite{lin2017_fpn} 将浅层细节和深层语义结合起来，后来许多小目标检测方法都沿用了这种多尺度思想。遥感场景中，SCRDet\cite{yang2019_scrdet} 等方法进一步结合注意力和旋转检测，以适应小尺度、密集分布目标。这类方法容易和现有检测器组合，但遇到极弱纹理目标时，特征仍可能被背景响应盖住。

另一类思路从尺度入手。SNIP\cite{singh2018_snip}、SNIPER\cite{singh2018_sniper} 和 AutoFocus\cite{najibi2019_autofocus} 的结果说明，训练样本选取、多尺度裁块和推理调度都会影响小目标表现。这类策略往往能改善候选区域分配，但实验流程随之变复杂。如果不同模型使用不同尺度策略，后续横向比较就必须格外谨慎。

还有一些工作不大改网络，而是调整训练目标和样本分布。Focal Loss\cite{lin2017_focalloss} 降低易分类样本权重，让模型更多关注困难样本；数据增强、困难样本重采样和损失重加权也常被用于小目标训练。它们实现成本相对低，但对数据质量和任务分布比较敏感。

最后一类是后处理和推理策略。小目标稀疏、背景又复杂时，置信度阈值、NMS 和融合策略会直接改变 Precision 与 Recall 的平衡。本文后续设置阈值消融，也是因为这些参数并非单纯的工程附属项。

已有综述\cite{wang2023_sod_review}也指出，遥感小目标检测还受到旋转变化、密集分布、复杂背景纹理和标注成本影响。因此，本文分析模型差异时不只讨论网络结构，也把输入尺度、样本分布和后处理参数一并纳入。

\subsection{遥感船舶检测研究现状}

船舶检测是遥感目标检测中应用较早、需求也比较明确的一类任务。随着公开数据集和开源框架增多，研究重心已经从手工特征逐步转到深度学习检测器。Zhao 等人\cite{zhao2024_shipsurvey} 对光学遥感船舶检测做过系统梳理，其中一个结论是：这类任务既继承通用检测框架，又必须处理近岸干扰、尺度变化、旋转形态和复杂海况等遥感特有问题。

LEVIR-Ship 数据集\cite{chen2022_tinyship_drenet} 聚焦中分辨率遥感图像中的微小船舶，为不同方法提供了同一数据基础。DRENet 也是围绕该数据集提出的专项方法，其退化重建增强机制主要服务于训练阶段，希望让网络对弱纹理船舶目标更敏感。本文选择它作为主线模型之一，正是因为它和本文场景最贴近。

工程工具链的成熟也改变了这类课题的组织方式。MMDetection\cite{mmdet2018_openmmlab} 和 Ultralytics\cite{ultralytics2024_docs} 已经提供了训练、测试和部署所需的基础设施，使不同模型之间的复现和比较更可操作。相应地，论文不能只报告一个模型分数，还需要说明评测口径、结果来源和系统接入方式。

\begin{table}[htbp]
  \centering
  \caption{本文相关研究方向与代表路线}
  \label{tab:ch2_research_map}
  \begin{tabular}{lll}
    \toprule
    研究方向 & 代表路线 & 本文定位 \\
    \midrule
    通用检测 & Faster R-CNN / FCOS / YOLO / DETR & 作为跨范式对比基础 \\
    小目标检测 & 多尺度融合、损失重加权、增强策略 & 用于分析性能差异来源 \\
    遥感船舶检测 & LEVIR-Ship + 专项方法（如 DRENet） & 聚焦微小船舶场景 \\
    工程化实现 & MMDetection / Ultralytics / 服务化封装 & 支撑系统落地与演示 \\
    \bottomrule
  \end{tabular}
\end{table}

\subsection{现有工作的不足与本文切入点}

从上述研究看，遥感微小船舶检测已有不少方法积累，但落到本文这样的多模型比较和系统实现任务时，仍有几个问题需要单独处理。

其一，实验协议常不完全一致。数据划分、输入尺度、阈值和评测脚本稍有不同，结果就很难直接并排解释。其二，一些研究对训练配置、日志记录和部署路径交代不够，复现时容易卡在工程细节上。其三，误差分析经常停留在总体指标，复杂海况、近岸背景和弱纹理目标造成的误检漏检没有被充分展开。其四，离线实验到可演示系统之间还有距离，模型结果如何进入任务管理、可视化和历史回看流程，很多论文并没有写清楚。

针对这些问题，本文以统一协议下的多模型比较为主线，选取 DRENet、FCOS 和 YOLO26 作为三个参照，并结合指标、样例和消融结果讨论它们的适用边界。随后再把可用推理能力接入轻量级检测系统，使算法结果能够被提交、展示和回看。

\subsection{小结}

本章围绕通用目标检测、小目标检测和遥感船舶检测梳理了相关研究，并说明本文为什么需要同时关注评测口径、复现记录、误差解释和系统接入。下一章将回到本文自己的实验设置，说明数据集、评价指标和模型方案。
